\documentclass[11pt]{article}

\usepackage[T1]{fontenc}
\usepackage[utf8]{inputenc}
\usepackage{lmodern}
\usepackage[margin=1in]{geometry}
\usepackage{microtype}
\usepackage{enumitem}
\usepackage[round,authoryear]{natbib}
\usepackage[hidelinks]{hyperref}

\title{Axodus: A Governance-Oriented Infrastructure Model for Human--AI
Coordination, Proof-of-Knowledge Participation, and Modular Digital
Institutions}
\author{Mauricio Z. Filho\\
\small Independent Researcher, Axodus\\
\small Correspondence: [public contact to be added]}
\date{Draft v0.1}

\begin{document}
\maketitle

\begin{center}
\fbox{\parbox{0.88\linewidth}{\centering
\textbf{Status:} Human review required; not ready for publication.\\
\textbf{Document type:} Conceptual architecture and research-agenda paper.\\
\textbf{License:} CC BY 4.0.}}
\end{center}

\begin{abstract}
Digital institutions increasingly combine human decision makers, automated
services, decentralized infrastructure, financial mechanisms, and public
documentation. Their technical components may be composable while their
authority, accountability, learning requirements, and risk controls remain
fragmented. This paper presents Axodus as a conceptual, prototype-stage
socio-technical architecture for examining that coordination problem. The
model organizes explicit governance rules, modular institutional components,
evidence-linked decisions, human-supervised AI assistance, educational
participation, and risk boundaries within a common research frame. It also
proposes separate L0--L5 and D0--D5 axes for documentation and development
maturity, without defining or assigning individual levels in this draft.
``Proof of Knowledge'' denotes a proposed educational participation mechanism,
not a cryptographic proof protocol. The contribution is architectural and
methodological: no operational deployment, user adoption, empirical outcome,
security property, or financial performance is claimed. The paper derives
research questions and an evaluation program covering traceability,
participation, human--AI reliance, maturity assessment, and risk governance.
Its purpose is to make a broad institutional design inspectable and
falsifiable before implementation claims are made.
\end{abstract}

\noindent\textbf{Keywords:} socio-technical systems; digital institutions;
human--AI coordination; governance; accountability; decentralized
organizations; educational participation; maturity models

\section{Introduction}

Digital ecosystems can execute transactions, coordinate distributed
contributors, and produce large volumes of documentation. These capabilities
do not by themselves determine who may decide, what evidence a decision
requires, how an affected party can challenge it, or who remains accountable
when automation contributes to an outcome. Socio-technical systems research
therefore treats organizational and technical design as interdependent rather
than assuming that one can be derived from the other
\citep{baxter2011sociotechnical}.

The problem becomes sharper when an institution combines decentralized
governance and AI assistance. Research on decentralized autonomous
organizations shows that self-executing rules remain embedded in social
processes for interpretation, change, and coordination
\citep{hassan2021dao,hsieh2018bitcoin}. Human--AI interaction research
likewise shows that useful automation requires expectation setting,
correction, and feedback mechanisms rather than an undifferentiated transfer
of authority \citep{amershi2019guidelines}.

This paper introduces Axodus as a conceptual architecture and research agenda
for these intersecting concerns. The project originates from a multi-year
independent research and design effort. The present contribution is limited to
a model: it does not report a functioning platform, production deployment,
users, evaluation results, or validated controls.

The paper makes three bounded contributions. First, it presents a layered
architecture connecting institutional rules, evidence, modular components,
educational participation, AI assistance, and risk boundaries. Second, it
separates documentation maturity from development maturity to reduce the risk
that extensive documentation is mistaken for operational readiness. Third, it
defines an evaluation agenda through which the proposed architecture could be
tested, rejected, or revised.

\section{Background and Related Work}

\subsection{Socio-technical and digital institutions}

A socio-technical perspective asks how roles, incentives, work practices,
information, and technology jointly shape outcomes. Baxter and Sommerville
argue for integrating organizational concerns with systems engineering rather
than treating them as a late adoption problem
\citep{baxter2011sociotechnical}. Axodus adopts this perspective by modeling
governance and accountability as architectural elements, not external policy
documents attached after implementation.

The term \emph{digital institution} is used here for an organized system of
roles, rules, records, and technical capabilities through which participants
pursue collective activity. It does not require that every rule be executable
or that every interaction occur on a blockchain.

\subsection{Decentralized governance and finance}

Hassan and De Filippi describe a DAO as a blockchain-based system in which
people coordinate through decentralized governance and self-executing rules
\citep{hassan2021dao}. Hsieh et al.\ distinguish machine consensus from the
social consensus required to alter protocols and organizational arrangements
\citep{hsieh2018bitcoin}. These distinctions caution against equating
technical distribution with accountable governance.

DeFi demonstrates how composable smart contracts can create open financial
services while introducing dependencies and technical, economic, and
governance risks \citep{schar2021defi}. Axodus does not claim to operate
financial services. DeFi is instead a motivating domain in which unclear
authority, uninformed participation, and weakly explained risk boundaries can
have consequential effects.

\subsection{Human--AI interaction and governance}

Guidelines for human--AI interaction emphasize communicating system
capabilities, supporting correction, and learning from behavior over time
\citep{amershi2019guidelines}. At the organizational level, the NIST AI Risk
Management Framework treats governance, mapping, measurement, and management
as connected risk functions \citep{tabassi2023airmf}. These sources motivate
an architecture in which AI outputs are attributable inputs to decisions
rather than sovereign decisions.

Constitutional AI uses written principles to guide model behavior through
supervised learning and AI feedback \citep{bai2022constitutional}. Axodus uses
\emph{constitutional} differently: it refers to institutional rules about
authority, prohibitions, amendments, review, and appeal. The proposed model
does not claim to train or align an AI model through Constitutional AI.

\subsection{Credentials and participation}

The W3C Verifiable Credentials Data Model separates the cryptographic
verification of a credential from validation of whether its claims are fit for
a verifier's purpose \citep{w3c2025vcdm}. This is central to the educational
component proposed here. A technically verifiable record would not establish
that an assessment is fair, that learning occurred, or that a governance
restriction is legitimate.

\section{Problem Statement}

The architecture addresses a design problem rather than an established
empirical effect:

\begin{quote}
How can a modular digital institution connect explicit authority,
evidence-linked accountability, educational participation, human-supervised
AI assistance, and risk boundaries without conflating documentation with
implementation or automation with legitimate governance?
\end{quote}

Five tensions structure the problem:

\begin{enumerate}[leftmargin=*]
\item \textbf{Modularity versus coherence.} Independent components can limit
coupling, but shared decisions require common identity, evidence, and review
rules.
\item \textbf{Participation versus readiness.} Open participation may support
legitimacy, while consequential decisions may require demonstrated
understanding.
\item \textbf{Automation versus accountability.} AI assistance can reduce
analytical load, but responsibility can become unclear when its outputs
influence action.
\item \textbf{Transparency versus privacy and security.} Traceable decisions
need records, while indiscriminate disclosure can harm participants or expose
systems.
\item \textbf{Documentation versus operational truth.} A detailed
specification can improve review while creating a misleading impression of
completion.
\end{enumerate}

Axodus is proposed as one architecture for making these tensions explicit. It
is not presented as their solution.

\section{Conceptual Architecture}

The model comprises six logical layers. They describe responsibilities and
interfaces, not an implementation topology.

\subsection{Constitutional governance}

The governance layer states roles, decision rights, prohibitions, amendment
rules, escalation paths, and appeal mechanisms. Its purpose is to make
higher-order authority visible. A valid action would need both technical
authorization and institutional justification; neither alone would establish
legitimacy.

\subsection{Evidence and accountability}

The evidence layer links proposals, supporting material, declared conflicts,
AI-generated analyses, decisions, responsible roles, and subsequent reviews.
The design objective is reconstructability: a reviewer should be able to ask
what was known, who was authorized, what assistance was used, and why an
outcome followed. The proposal does not assume that more data automatically
produces accountability.

\subsection{Modular institutional components}

Bounded components represent institutional capabilities such as learning,
proposal review, risk analysis, or record keeping. Each component should
declare its purpose, inputs, outputs, authority, dependencies, evidence
requirements, and failure behavior. No definitive inventory or implementation
status is asserted in this draft.

\subsection{Participation and learning}

An Academy component and Proof of Knowledge mechanism are proposed as
interfaces for learning and assessment. Their purpose would be to help
participants understand domain-specific risks and governance duties before
performing selected consequential actions. They must not become a general
social ranking system.

\subsection{Human--AI assistance}

AI services may support bounded tasks such as document comparison, proposal
summarization, inconsistency detection, risk flagging, and preparation of
review material. Each use should declare the task, source material, output,
uncertainty where available, human reviewer, and disposition. Humans retain
decision authority and responsibility.

\subsection{Risk and sustainability boundaries}

The risk layer expresses constraints on authority, exposure, timing,
conflicts, and escalation. In a future treasury context, the same layer could
require auditable rationales and approvals. This is a design requirement only:
the paper claims no working treasury, enforced limit, security property, or
financial outcome.

\section{Modularity and Maturity}

Axodus proposes two distinct maturity axes. L0--L5 represents documentation
maturity; D0--D5 represents development maturity. This paper deliberately does
not define the individual levels, assign scores, or report component status.
Those actions require observable criteria and evidence not available in the
current public repository.

The separation is intended to prevent three category errors:

\begin{itemize}[leftmargin=*]
\item treating a comprehensive specification as implemented behavior;
\item treating an implemented prototype as validated institutional
performance;
\item averaging unlike forms of maturity into one reassuring score.
\end{itemize}

A future maturity specification should define required artifacts, admissible
evidence, reviewers, expiration rules, and disagreement handling for every
level. Inter-rater reliability and resistance to self-assessment inflation
should be evaluated before the axes are used in public reporting.

Modularity is similarly a hypothesis rather than an assured property.
Component boundaries may support bounded change and review, but they can also
move complexity into interfaces. Evaluation must therefore examine dependency
visibility, cross-component failure, and governance consistency.

\section{Proof of Knowledge and Educational Participation}

\noindent\textbf{Terminology note.} \emph{Proof of Knowledge} in this paper
means a proposed educational participation mechanism. It is not a
cryptographic proof of knowledge, proof-of-work, proof-of-stake, or a deployed
credential protocol.

The mechanism would connect a defined learning objective, an assessment, a
reviewable result, and eligibility for a bounded activity. A possible flow is:

\begin{enumerate}[leftmargin=*]
\item publish the knowledge and risk objectives for an activity;
\item provide accessible learning paths;
\item assess understanding using methods appropriate to the objective;
\item issue a time-bounded result with an appeal path;
\item apply the result only to the specified activity;
\item reassess validity, disparate effects, and continued necessity.
\end{enumerate}

The W3C credential model could inform portable representations, but
verification of an issuer and credential does not validate educational quality
or justify authorization \citep{w3c2025vcdm}. The governance problem therefore
includes who defines curricula, who assesses assessors, how accommodations are
provided, and how a participant can contest an outcome.

The central hypothesis is that education-linked readiness may improve the
quality of participation in selected high-consequence contexts. Competing
hypotheses are equally important: the mechanism may reduce participation,
favor already advantaged groups, encourage test gaming, centralize curriculum
power, or create privacy risks. Proof of Knowledge should be rejected or
redesigned if those harms cannot be controlled.

\section{Governance, Accountability, and Human Oversight}

The proposed decision record has seven conceptual fields: decision scope,
authorized roles, evidence considered, conflicts disclosed, automated
assistance used, decision and rationale, and review or appeal path. This is a
minimum accountability model, not a database schema.

AI assistance is governed by task-specific boundaries. Low-consequence
drafting may permit lightweight review; risk classification or proposal
analysis would require stronger provenance, independent checking, and explicit
human disposition. An AI output cannot approve its own use, resolve an appeal
about itself, or become the final bearer of responsibility.

The model should support contestability. Participants affected by a decision
need a way to identify applicable rules, challenge evidence, correct records,
and request review by an appropriately independent role. Immutable records,
where used, must not be treated as immutable interpretations.

These requirements align with risk-oriented governance and human--AI
interaction guidance \citep{tabassi2023airmf,amershi2019guidelines}, but their
effectiveness in Axodus remains untested.

\section{Risk Boundaries and Treasury Transparency}

Risk boundaries are proposed at four levels:

\begin{itemize}[leftmargin=*]
\item \textbf{action boundaries:} what an actor or agent may do;
\item \textbf{exposure boundaries:} resources or values that may be placed at
risk;
\item \textbf{process boundaries:} evidence, approvals, waiting periods, and
conflicts;
\item \textbf{escalation boundaries:} conditions requiring suspension or
review.
\end{itemize}

In financial settings, on-chain visibility can reveal transactions while
leaving assumptions, conflicts, and rationales opaque. DeFi research also
shows that composability can transmit dependencies across protocols
\citep{schar2021defi}. A future Axodus treasury model would therefore need
both transaction records and decision context.

No treasury architecture, threshold, asset policy, token mechanism, or
financial projection is disclosed or claimed here. Before any live-value
experiment, the model would require threat modeling, simulation, independent
security review, legal analysis, incident procedures, and explicit human
authorization.

\section{Prototype Status}

Axodus is presented in this paper as a conceptual, prototype-stage model. The
public repository presently supports editorial artifacts and research
documentation. It does not supply evidence of an integrated implementation,
production operation, user adoption, empirical performance, or effective
controls.

The absence of those claims is substantive. Architectural consistency can be
reviewed before implementation, but it cannot substitute for observed
behavior. Future versions should maintain a component-by-component evidence
table distinguishing designed, prototyped, tested, deployed, and independently
evaluated states.

\section{Evaluation Plan}

The architecture yields five initial research questions.

\paragraph{RQ1---Traceability.} Does the proposed evidence model help
independent reviewers reconstruct decisions more accurately and efficiently
than a document-only baseline? A study could use matched governance scenarios,
blinded reviewers, reconstruction accuracy, time, disagreement, and
missing-evidence detection.

\paragraph{RQ2---Educational participation.} Does Proof of Knowledge improve
domain-relevant understanding without unacceptable exclusion? A staged study
should measure learning, retention, decision reasoning, completion,
accessibility, subgroup effects, appeals, and test gaming. Token or financial
stakes should not be introduced in an initial study.

\paragraph{RQ3---Human--AI reliance.} For bounded proposal-analysis tasks, how
do provenance and mandatory human disposition affect error detection,
calibrated reliance, override behavior, and review quality? Experiments should
include planted AI errors and ambiguous cases.

\paragraph{RQ4---Maturity reliability.} Can independent reviewers apply future
L and D criteria consistently from the same evidence? Evaluation should
measure inter-rater agreement, reasons for disagreement, and susceptibility to
strategic self-rating.

\paragraph{RQ5---Risk governance.} Do explicit action, exposure, process, and
escalation boundaries reduce unauthorized or poorly evidenced decisions in
tabletop simulations? Scenarios should include conflicts of interest, urgent
actions, incomplete evidence, agent error, and cross-component failure.

The sequence should move from document inspection to tabletop simulation,
non-financial prototypes, controlled user studies, and only then bounded field
pilots. Each stage needs preregistered success and stop criteria appropriate to
its risk.

\section{Limitations and Threats to Validity}

This paper is broad and conceptual. It may combine concerns that require
separate theories, methods, and papers. The architecture can appear coherent
because it has not yet encountered implementation constraints or institutional
conflict.

The public evidence base for Axodus is currently limited to design
documentation. There are no results with which to assess feasibility,
usability, governance legitimacy, security, or social impact. The model is
also authored from a single-project perspective and may encode unexamined
assumptions about authority, education, and accountability.

Proof of Knowledge presents particular risks of exclusion, surveillance,
curriculum capture, and false confidence. Human oversight can become nominal
when reviewers lack time or expertise. Audit trails can create privacy and
security harms. Modularity can hide systemic coupling. Explicit rules can
formalize unjust arrangements rather than correct them.

Web3-facing governance designs may also encounter regulatory uncertainty,
operational fragility, and jurisdiction-specific compliance burdens. AI
assistance may be unreliable, difficult to calibrate, or difficult to audit
under realistic workload and time pressure.

The paper does not provide legal, financial, security, or investment advice.
Its terminology may not map cleanly onto every jurisdiction or research
community.

\section{Research Agenda}

Near-term work should:

\begin{enumerate}[leftmargin=*]
\item formalize observable L and D maturity criteria;
\item specify the minimal decision and evidence model;
\item design a non-financial Proof of Knowledge prototype with accessibility
and appeal requirements;
\item define bounded AI-assistance tasks and provenance records;
\item create governance and failure scenarios for tabletop evaluation;
\item recruit independent reviewers from governance, HCI, education,
security, and decentralized-systems backgrounds;
\item publish negative findings and design changes alongside positive results.
\end{enumerate}

Separate papers should address educational validity, human--AI governance, and
financial risk rather than treating the present architecture paper as evidence
in those domains.

\section{Conclusion}

Axodus is proposed as a conceptual architecture for connecting governance,
evidence, modular institutional components, educational participation,
human-supervised AI assistance, and risk boundaries. Its immediate value is
not demonstrated performance but inspectability: the model exposes claims,
interfaces, tensions, and evaluation questions before operational maturity is
asserted.

The architecture will become a credible contribution only if its constructs
are specified, independently criticized, implemented in bounded prototypes,
and evaluated against alternatives. This draft establishes a research agenda
for that process. It is ready for human review, not for public publication.

\bibliographystyle{plainnat}
\bibliography{references}

\end{document}
